\paragraph{What the ontology explains for scientific users.}
The useful association is between a scientific requirement and its recorded implementation, rather than between a domain label and a presumed bottleneck. Jacobi's update-norm diagnostic connects a PDE requirement to two reduction policies: a native block/shuffle tree and Julia per-point floating atomics. Their different counted FP32 numerators reflect reduction lowering and the counter boundary, not additional field updates. STREAM dot instead has equal counted work across versions but substantially different cold-replay traffic, while Triad's integrated intensities are nearly equal. These paired observations motivate case-specific investigation; they do not establish a universal abstraction penalty. Memory placement supplies another necessary distinction: CPU-resident mapped data require a host/PCIe view, rather than an HBM-only ceiling. The ontology therefore links scientific stage, accuracy contract, reduction, precision, memory tier and backend to qualified observations. This helps users identify which implementation choices require further measurement while preserving missing data and unresolved mechanisms. Scientific-family coverage is descriptive; neither partial stage samples nor these three paired families justify a corpus-wide intensity--speedup correlation or a domain-wide framework recommendation.
